% =============================================================
%  Term of Reference — Proyek Akhir Multimodal
%  IF25-40304 · Pembelajaran Mesin Multimodal · ITERA
%  Pengganti UAS untuk paruh kedua semester (Pertemuan 09--16)
% =============================================================
\documentclass[12pt, a4paper]{article}

\usepackage{tugas-style}

% ── Identitas dokumen ───────────────────────────────────────
\renewcommand{\doctitle}{Term of Reference Proyek Akhir}
\renewcommand{\duedate}{\color{ctpBlue}\faCalendarCheck\ \ Rentang proyek: Pertemuan 09--16 \textbar\ Deliverable akhir pada UAS \faCalendarCheck}

% Header: dokumen ini bukan "Tugas Individu", jadi label disesuaikan
\lhead{{\footnotesize\color{ctpOverlay}\textsc{Proyek Akhir Multimodal}}}
\rhead{{\footnotesize\color{ctpOverlay}\textsc{\coursenumber\ --\ \coursetitle}}}

\begin{document}

% ── HEADER ──────────────────────────────────────────────────
\makeassignmentheader

\hrule width0.3\textwidth
\revisioninfo{%
  \vhEntry{0.1}{24 September 2026}{I.W.W.}{Draf inisial Term of Reference.}
}
\hrule
\vskip .3in

\tableofcontents
\clearpage

% =============================================================
\section{Ringkasan Proyek}\label{sec:ringkasan}

Proyek Akhir Multimodal adalah penilaian paruh kedua semester yang
\textbf{menggantikan UAS tertulis}. Mahasiswa bekerja dalam kelompok untuk
membangun, mengevaluasi, dan mempertanggungjawabkan sebuah sistem
\textbf{pencarian lintas modal} (\emph{cross-modal retrieval}) pada data budaya
Indonesia, lalu mempresentasikan hasilnya pada sesi UAS.

\renewcommand{\arraystretch}{1.35}
\begin{center}\small
\begin{tabular}{|L{3.9cm}|L{10.5cm}|}
  \hline
  \thead{\color{ctpBlue} Aspek} & \thead{\color{ctpBlue}Ketentuan} \\
  \hline
  Bentuk & Kelompok, \textbf{5 mahasiswa} per kelompok (10 kelompok per kelas). \\
  \hline
  Data & SEA-VL \emph{Indonesian subset} --- 2.763 gambar budaya Indonesia dengan caption
         Indonesia dan Inggris. \emph{Split} latih dan validasi dibagikan; \emph{test set}
         (416 gambar) ditahan dosen. \\
  \hline
  Tugas & Pencarian dua arah gambar$\leftrightarrow$teks dengan kueri berbahasa Indonesia
         dan klasifikasi kategori budaya secara \emph{zero-shot}. \\
  \hline
  Metrik & \emph{Recall@1/5/10} dan \emph{mAP@K}. \\
  \hline
  Aturan utama & Encoder pra-latih \textbf{dibekukan}. Tidak melatih dari nol, tidak
         mengunduh dataset baru. \\
  \hline
  Jangka waktu & Pertemuan 09--16 dengan lima titik periksa (M0--M4). \\
  \hline
  Bobot terhadap nilai akhir & \textbf{25\%} titik periksa + \textbf{25\%} deliverable akhir. \\
  \hline
\end{tabular}
\end{center}

% =============================================================
\section{Tujuan Pembelajaran}\label{sec:objectives}

Setelah menyelesaikan proyek ini, mahasiswa mampu:
\begin{objectives}
  \item Merancang dan melaksanakan proyek multimodal \emph{end-to-end}, mulai dari
        rumusan masalah, pemilihan encoder, hingga pelaporan hasil.
  \item Menerapkan ketiga pilar mata kuliah --- \emph{representasi}, \emph{penyejajaran},
        dan \emph{fusi} --- pada data budaya nyata dan mengukur kontribusi tiap pilar.
  \item Mengevaluasi sistem pencarian lintas modal dengan protokol yang bebas kebocoran,
        termasuk memisahkan \emph{split} acak dari \emph{split} independen-penutur.
  \item Menjelaskan \emph{kesenjangan lintas bahasa} secara kuantitatif dan menilai
        implikasi etisnya bagi pengguna berbahasa Indonesia.
  \item Bekerja dalam tim dengan pembagian peran yang eksplisit dan kontribusi yang
        dapat diverifikasi secara individual.
  \item Mempertanggungjawabkan keputusan teknis secara lisan dan tertulis.
\end{objectives}

% =============================================================
\section{Catatan dan Batasan}\label{sec:notes}
\vspace{-.1cm}

\begin{minipage}{\linewidth}
\begin{bclogo}[couleur=ctpBase, arrondi=0.1, logo=\bccrayon, ombre=false, couleurBord=ctpSurface, epBord=0.6]{Perhatikan hal berikut:}
  \begin{coloredPen}
    \item Proyek bersifat \textbf{kelompok} (5 mahasiswa). Setiap anggota \textbf{wajib}
          memegang satu peran dengan artefak yang dapat diperiksa (lihat
          Bagian~\ref{sec:peran}).
    \item Encoder pra-latih \textbf{wajib dibekukan}. Melatih ulang atau memodifikasi
          bobot encoder dinilai sebagai kesalahan prosedur.
    \item \textbf{Tidak diperkenankan} mengunduh, membuat, atau menambal dataset baru,
          maupun menambahkan gambar dari sumber luar.
    \item \emph{Test set} bersifat \textbf{privat} dan baru dibuka pada Pertemuan 13.
          Segala bentuk upaya mengakses atau menebaknya sebelum tenggat dinilai sebagai
          pelanggaran integritas akademik.
    \item Setiap eksperimen harus selesai \textbf{di bawah 10 menit pada CPU}. GPU
          diperbolehkan bila tersedia, bukan syarat.
    \item Seluruh hasil angka yang dilaporkan harus berasal dari \textbf{eksekusi nyata}.
          Angka yang dikarang atau disalin dari makalah pihak lain tanpa menjalankan
          ulang akan digugurkan.
    \item Data berlisensi \textbf{CC BY-SA 4.0}. Cantumkan atribusi pada setiap laporan,
          \emph{notebook}, dan slide.
  \end{coloredPen}
\end{bclogo}
\end{minipage}

\begin{minipage}{\linewidth}
\begin{bclogo}[couleur=ctpMantle, arrondi=0.1, logo=\bcinfo, ombre=false, couleurBord=ctpSurface, epBord=0.6]{Ekspektasi Hasil}
  Sebagai pembanding, inilah \emph{baseline} yang telah diukur dosen pada 500 gambar
  validasi dengan encoder beku \code{clip-ViT-B-32} dan \code{clip-ViT-B-32-multilingual-v1}.
  Tebakan acak hanya \textbf{0,20\%}.

  \vspace{2pt}
  \begin{center}\small
  \renewcommand{\arraystretch}{1.3}
  \begin{tabular}{|L{7.2cm}|c|c|c|}
    \hline
    \thead{\color{ctpBlue} Arah pencarian} & \thead{\color{ctpBlue}R@1} & \thead{\color{ctpBlue}R@5} & \thead{\color{ctpBlue}R@10} \\
    \hline
    gambar $\rightarrow$ teks \textbf{Inggris} (\emph{plafon pembanding}) & 35,6\% & 64,0\% & 72,6\% \\
    \hline
    gambar $\rightarrow$ teks \textbf{Indonesia} (\emph{baseline}) & \textbf{15,6\%} & \textbf{30,8\%} & \textbf{37,4\%} \\
    \hline
    teks \textbf{Inggris} $\rightarrow$ gambar & 29,6\% & 56,6\% & 66,2\% \\
    \hline
    teks \textbf{Indonesia} $\rightarrow$ gambar & \textbf{14,6\%} & \textbf{30,4\%} & \textbf{39,8\%} \\
    \hline
  \end{tabular}
  \end{center}

  Perhatikan \textbf{jurang 20 poin persen} pada \emph{Recall@1} antara kueri berbahasa
  Inggris dan Indonesia. Menaikkan angka Indonesia itulah inti tantangan proyek ini ---
  bukan sekadar memilih model yang paling besar.
\end{bclogo}
\end{minipage}

% =============================================================
\section{Peran Anggota Kelompok}\label{sec:peran}

Setiap anggota memegang \textbf{satu} peran. Peran tidak boleh dirangkap, dan setiap
peran menghasilkan artefak tersendiri yang dinilai pada titik periksa.

\renewcommand{\arraystretch}{1.35}
\begin{center}\small
\begin{tabular}{|c|L{3.6cm}|L{8.6cm}|}
  \hline
  \thead{\color{ctpBlue}\#} & \thead{\color{ctpBlue}Peran} & \thead{\color{ctpBlue}Artefak yang diperiksa} \\
  \hline
  1 & Data dan Kurasi & Kode penyaringan data, laporan kualitas data, penanganan duplikat dan lisensi. \\
  \hline
  2 & Baseline dan Pipeline & Kode \emph{end-to-end} yang berjalan, skrip reproduksi, tabel hasil awal. \\
  \hline
  3 & Eksperimen Penyejajaran/Fusi & Matriks eksperimen, kode ablasi, kurva hasil. \\
  \hline
  4 & Evaluasi dan Metrik & Protokol evaluasi, tabel \emph{Recall@K}/\emph{mAP}, analisis kuantitatif. \\
  \hline
  5 & Integrasi, Etika, dan Koordinasi & Laporan, slide presentasi, pernyataan limitasi dan etika, log kontribusi. \\
  \hline
\end{tabular}
\end{center}

\begin{callout}[ctpBlue]
\textbf{\color{ctpBlue}Akuntabilitas Individu}
\begin{itemize}
  \item Setiap anggota mengisi \textbf{log mingguan} (satu baris per minggu) yang
        mencatat pekerjaan, bukti, dan kendala.
  \item Pada sesi tanya jawab, pertanyaan diarahkan \textbf{langsung kepada pemegang
        peran terkait}, bukan hanya kepada ketua kelompok.
  \item Faktor kontribusi dari penilaian sebaya dikalikan pada komponen kontribusi
        individu, dengan rentang sempit (0,85--1,15) agar tidak menjadi alat balas dendam.
\end{itemize}
\end{callout}

% =============================================================
\section{Perangkat Lunak dan Peralatan yang Diperlukan}\label{sec:requiredsw}

\begin{itemize}[itemsep=2pt,parsep=0pt,topsep=2pt,partopsep=2pt]
  \item[\color{ctpBlue}\faPython] \textbf{Python} 3.10 atau lebih baru.
  \item[\color{ctpBlue}\faFileCode] \textbf{Paket wajib:} \code{numpy}, \code{torch},
        \code{transformers}, \code{pillow}, \code{pandas}.
  \item[\color{ctpBlue}\faFileCode] \textbf{Alternatif:} \code{open\_clip} sebagai pengganti
        \code{transformers} untuk encoder CLIP.
  \item[\color{ctpBlue}\faDownload] \textbf{Data:} bundel SEA-VL Indonesia diunduh dari LMS
        (berisi \code{images/}, \code{train.csv}, \code{val.csv}, \code{ATTRIBUTION.md}).
        \emph{Test set} tidak disertakan.
  \item[\color{ctpBlue}\faFileCode] \textbf{Perangkat:} CPU memadai. Pengukuran dosen
        menunjukkan penyematan seluruh 3.096 gambar memerlukan waktu sekitar 80 detik
        pada perangkat keras kelas laptop.
\end{itemize}

\begin{codewindow}{setup.sh}
\begin{lstlisting}[style=pycode]
pip install numpy torch transformers pillow pandas
\end{lstlisting}
\end{codewindow}

\begin{minipage}{\linewidth}
\begin{bclogo}[couleur=ctpMantle, arrondi=0.1, logo=\bcinfo, ombre=false, couleurBord=ctpSurface, epBord=0.6]{Batasan Cakupan}
  Proyek ini \textbf{bukan} proyek pelatihan model dari nol. Fokusnya adalah
  \emph{rekayasa penyejajaran dan fusi} di atas representasi beku: bagaimana menyusun
  proyeksi, prompt, dan mekanisme pemeringkatan ulang sehingga gambar dan teks
  berbahasa Indonesia bertemu pada tempat yang benar di ruang vektor.
\end{bclogo}
\end{minipage}

% =============================================================
\section{Pernyataan Masalah}\label{sec:intro}

Indonesia memiliki warisan budaya visual yang sangat kaya: rumah adat, tarian, tekstil,
kuliner, hingga benda cagar budaya. Sebagian besar dokumentasinya tersebar sebagai foto
dengan deskripsi pendek. Seorang pengguna ingin mencari ``\emph{tari topeng dari Cirebon
dengan penari berkostum merah}'' dan mengharapkan sistem mengembalikan foto yang tepat
--- tanpa kata kunci yang persis sama.

Persoalannya adalah \textbf{penyejajaran lintas modal}: memastikan sebuah gambar dan
sebuah kalimat yang membicarakan hal yang sama berada pada titik yang berdekatan di
ruang vektor. Sesi Pertemuan 07 dan 09 menunjukkan bahwa model \emph{dual-encoder}
seperti CLIP melakukannya melalui \emph{cross-attention} dan tujuan kontrastif, tanpa
pernah mempertemukan lapisan kedua modalitas.

Namun ada twist yang membuat persoalan ini menjadi relevan dan belum selesai. Sistem
semacam ini dirancang untuk kueri berbahasa Inggris. Ketika pengguna mengetik dalam
bahasa Indonesia, kinerjanya \textbf{anjlok sekitar 20 poin persen}. Analisis Anda
harus menjelaskan mengapa hal itu terjadi dan apa yang dapat dilakukan untuk
mempersempit jurang tersebut.

% =============================================================
\section{Persyaratan}\label{sec:requirements}

Proyek dilaksanakan dalam lima titik periksa. Setiap titik periksa dinilai
\textbf{5 poin} dari 100 poin mata kuliah. Titik periksa berikutnya tidak dapat
dinilai apabila titik periksa sebelumnya tidak dikumpulkan.

\renewcommand{\arraystretch}{1.3}
\begin{center}\small
\begin{tabular}{|c|L{1.5cm}|L{5.3cm}|L{6.0cm}|}
  \hline
  \thead{\color{ctpBlue}Kode} & \thead{\color{ctpBlue}Pertemuan} & \thead{\color{ctpBlue}Fokus} & \thead{\color{ctpBlue}Yang dikumpulkan} \\
  \hline
  M0 & 09 & Proposal dan reproduksi \emph{baseline} & 1 halaman proposal + angka \emph{baseline} hasil reproduksi mandiri. \\
  \hline
  M1 & 10 & \emph{Baseline} pencarian berjalan & Kode yang berjalan \emph{end-to-end} + tabel \emph{Recall@K} pada \emph{val}. \\
  \hline
  M2 & 11 & Perbandingan strategi penyejajaran/fusi & Laporan singkat $\geq$2 pendekatan + tabel perbandingan. \\
  \hline
  M3 & 12 & Uji kegagalan dan etika & Hasil uji kegagalan + pernyataan limitasi, etika, dan atribusi. \\
  \hline
  M4 & 13 & Protokol evaluasi dan \emph{test set} & Protokol evaluasi + pengajuan prediksi pada \emph{test set} privat. \\
  \hline
\end{tabular}
\end{center}

% ─────────────────────────────────────────────────────────────
\subsection{M0 --- Proposal dan Reproduksi Baseline}

Titik periksa ini membuka proyek sekaligus mengukur kesiapan teknis kelompok.

\begin{enumerate}
  \item Tulis proposal \textbf{satu halaman} yang memuat: rumusan masalah, pertanyaan
        penelitian, pasangan modalitas yang digarap, metrik yang dipakai, dan rencana
        lima sumbu eksperimen kelompok.
  \item Jalankan kembali perhitungan \emph{baseline} dari lembar kerja Pertemuan 09 dan
        laporkan angka Anda sendiri untuk \emph{InfoNCE} dan \emph{Recall@1}.
  \item Sebutkan lima peran kelompok beserta nama anggotanya.
\end{enumerate}

% ─────────────────────────────────────────────────────────────
\subsection{M1 --- Baseline Pencarian Lintas Modal}

\begin{enumerate}
  \item Bangun pipeline \emph{end-to-end}: baca \code{train.csv}/\code{val.csv}, sematkan
        gambar dan caption, hitung kesamaan kosinus, kembalikan peringkat.
  \item Laporkan \emph{Recall@1/5/10} dan \emph{mAP@K} pada \emph{split} validasi untuk
        \textbf{empat arah}: gambar ke teks Indonesia, gambar ke teks Inggris,
        teks Indonesia ke gambar, dan teks Inggris ke gambar.
  \item Tampilkan sepuluh contoh kueri berbahasa Indonesia beserta hasil teratasnya,
        termasuk setidaknya dua contoh \textbf{yang gagal}.
\end{enumerate}

\begin{codewindow}{latih\_baseline.py}
\begin{lstlisting}[style=pycode]
import pandas as pd, torch
from PIL import Image
from sentence_transformers import SentenceTransformer

img_model = SentenceTransformer("sentence-transformers/clip-ViT-B-32")
txt_model = SentenceTransformer("sentence-transformers/clip-ViT-B-32-multilingual-v1")

df = pd.read_csv("val.csv")
imgs  = [Image.open(f).convert("RGB") for f in df["file"]]
cap_id = list(df["caption_id"])
cap_en = list(df["caption_en"])

E  = img_model.encode(imgs,    batch_size=32, normalize_embeddings=True)
Ti = txt_model.encode(cap_id,  batch_size=64, normalize_embeddings=True)
Te = img_model.encode(cap_en,  batch_size=64, normalize_embeddings=True)

def recall(A, B, ks=(1, 5, 10)):
    rank = (A @ B.T).argsort(dim=1, descending=True)
    truth = torch.arange(len(A)).unsqueeze(1)
    return {k: float((rank[:, :k] == truth).any(1).float().mean()) for k in ks}

print("gambar -> teks-ID :", recall(torch.tensor(E), torch.tensor(Ti)))
print("teks-ID -> gambar :", recall(torch.tensor(Ti), torch.tensor(E)))
\end{lstlisting}
\end{codewindow}

% ─────────────────────────────────────────────────────────────
\subsection{M2 --- Perbandingan Strategi Penyejajaran dan Fusi}

Pilih dan bandingkan \textbf{sekurang-kurangnya dua} pendekatan untuk mempersempit
jurang lintas bahasa. Setiap kelompok \textbf{wajib menggarap satu sumbu yang berbeda}
agar kelas menghasilkan bukti yang beragam.

\begin{enumerate}
  \item Pilih satu sumbu dari daftar berikut, atau ajukan sumbu sendiri kepada dosen:
        \begin{alphalist}
          \item mengganti encoder multibahasa (misalnya model berbasis SigLIP atau
                \emph{multilingual CLIP} yang lebih baru);
          \item \emph{prompt ensemble} berbahasa Indonesia, termasuk terjemahan dan
                variasi templat;
          \item melatih proyeksi/adapter kecil di atas representasi beku;
          \item \emph{reranking} dengan \emph{cross-encoder} pada sejumlah kandidat teratas;
          \item penajaman data dan augmentasi terjemahan.
        \end{alphalist}
  \item Laporkan efek pendekatan pada keempat arah pencarian, dalam bentuk tabel
        sebelum--sesudah.
  \item Jelaskan biaya pendekatan Anda: waktu komputasi, jumlah parameter yang dilatih,
        dan kompleksitas penerapannya.
\end{enumerate}

% ─────────────────────────────────────────────────────────────
\subsection{M3 --- Uji Kegagalan, Etika, dan Atribusi}

\begin{enumerate}
  \item Rancang \textbf{uji kegagalan komposisional}: gunakan lima gambar dan kueri yang
        memakai kata serupa tetapi berbeda relasi, negasi, atau penanda budaya.
        Laporkan probabilitas setiap kandidat.
  \item Analisis apakah sistem lebih mudah menangani objek tunggal daripada relasi,
        dan hubungkan dengan keterbatasan arsitektur \emph{dual-encoder}.
  \item Tulis pernyataan \textbf{etika dan limitasi} (maks. satu halaman) yang membahas:
        \begin{alphalist}
          \item kesenjangan kinerja antara kueri Indonesia dan Inggris, serta siapa yang
                dirugikan olehnya;
          \item keterwakilan wilayah dan kelompok budaya di dalam data;
          \item kejujuran \emph{benchmark}: mengapa laporan tanpa \emph{split} yang jelas
                sulit dipercaya.
        \end{alphalist}
  \item Cantumkan atribusi CC BY-SA 4.0 dengan benar pada laporan dan kode.
\end{enumerate}

% ─────────────────────────────────────────────────────────────
\subsection{M4 --- Protokol Evaluasi dan Pengajuan Test Privat}

\begin{enumerate}
  \item Tulis \textbf{protokol evaluasi} yang menjelaskan: pembagian data, metrik,
        jumlah \emph{run}, cara menangani kelas yang jarang, dan bentuk tabel hasil akhir.
  \item Bandingkan \textbf{split acak} dengan \textbf{split yang dikelompokkan} (misalnya
        berdasarkan lokasi atau kategori) dan jelaskan selisihnya.
  \item Kirim prediksi kelompok Anda untuk \emph{test set} privat dalam format yang
        ditetapkan dosen. Satu pengajuan per kelompok.
\end{enumerate}

\begin{callout}[ctpPeach]
\textbf{\color{ctpPeach}Peringatan Metodologis}
\emph{Data leakage} adalah kesalahan paling umum pada proyek seperti ini. Periksa
kembali: apakah gambar duplikat masih tersisa; apakah caption yang sangat mirip muncul
di \emph{train} dan \emph{val} sekaligus; apakah Anda mengoptimalkan apa pun dengan
melihat \emph{val}. Laporkan temuan Anda, bukan hanya angka akhirnya.
\end{callout}

% ─────────────────────────────────────────────────────────────
\subsection{Klinik dan Gladi Bersih}

\begin{itemize}
  \item \textbf{Pertemuan 14 --- Klinik:} setiap kelompok membawa draf laporan dan
        melakukan \emph{peer review} terhadap satu kelompok lain menggunakan rubrik.
  \item \textbf{Pertemuan 15 --- Gladi bersih:} demo berjalan dengan batas waktu ketat.
        Umpan balik akhir diberikan sebelum penilaian.
\end{itemize}

% ─────────────────────────────────────────────────────────────
\subsection{Deliverable Akhir --- Sesi UAS (Pertemuan 16)}

\begin{enumerate}
  \item \textbf{Laporan teknis} (4--6 halaman) dengan urutan: rumusan masalah, data dan
        protokol, metode, hasil, ablasi, dan limitasi.
  \item \textbf{Kode} yang dapat dijalankan ulang, disertai berkas intruksi singkat
        (\code{README}) dan penetapan \emph{seed}.
  \item \textbf{Presentasi} 10 menit dengan demo langsung.
  \item \textbf{Sesi tanya jawab} 5 menit, dengan pertanyaan tertarget kepada pemegang
        peran.
\end{enumerate}

\begin{callout}[ctpGreen]
\textbf{\color{ctpGreen}Papan Peringkat dan Bonus Kompetisi}
\begin{itemize}
  \item Papan peringkat diperbarui pada setiap titik periksa, berdasarkan \emph{Recall@1}
        dan \emph{mAP@K} pada \emph{val}, serta menampilkan terpisah \textbf{peringkat
        efisiensi} (kinerja per biaya komputasi) agar kelompok tanpa GPU tetap dapat
        menang.
  \item Kategori yang dihargai selain peringkat teratas: \emph{analisis terbaik} dan
        \emph{paling berkembang}.
  \item Bonus kompetisi maksimum adalah \textbf{5\% dari total komponen proyek
        (setara 2,5 poin dari 100)}, tidak dapat menaikkan nilai melewati ambang batas
        tanpa peninjauan dosen, dan sama sekali tidak memengaruhi kelulusan.
\end{itemize}
\end{callout}

% =============================================================
\clearpage
\section{Kriteria Penilaian}\label{sec:evaluation}

\subsection{Bobot Mata Kuliah}

\renewcommand{\arraystretch}{1.45}
\begin{center}\small
\begin{tabular}{|L{10.4cm}|c|}
  \hline
  \thead{\color{ctpBlue} Komponen} & \thead{\color{ctpBlue}Bobot} \\
  \hline
  UTS (ujian tertulis, Pertemuan 08) & 25\% \\
  \hline
  Tugas pra-UTS (Pertemuan 02, 03, 04, 06) & 25\% \\
  \hline
  Progress proyek (titik periksa M0--M4, masing-masing 5\%) & 25\% \\
  \hline
  \quad --- Laporan dan kode akhir & 15\% \\
  \hline
  \quad --- Presentasi, tanya jawab, dan kontribusi individu & 10\% \\
  \hline
  \textbf{Total} & \textbf{100\%} \\
  \hline
\end{tabular}
\end{center}

\subsection{Rubrik Titik Periksa (M0--M4)}

Setiap titik periksa bernilai 5 poin. Nilai diberikan per kriteria dengan tingkat
0--3, lalu dikonversi ke bobot poin.

\renewcommand{\arraystretch}{1.35}
\begin{center}\small
\begin{tabular}{|L{3.4cm}|L{2.1cm}|L{2.0cm}|L{2.0cm}|L{2.0cm}|}
  \hline
  \thead{\color{ctpBlue}Kriteria} & \thead{\color{ctpBlue}Tingkat 0} & \thead{\color{ctpBlue}Tingkat 1} & \thead{\color{ctpBlue}Tingkat 2} & \thead{\color{ctpBlue}Tingkat 3} \\
  \hline
  Kelengkapan artefak & Tidak ada & Sebagian, tidak terdokumentasi & Lengkap & Lengkap dan rapi \\
  \hline
  Ketepatan teknis & Tidak berjalan & Berjalan dengan bantuan & Berjalan mandiri & Berjalan mandiri dan reproduksibel \\
  \hline
  Kualitas analisis & Tidak ada & Deskriptif & Analitis & Analitis dengan bukti kuantitatif \\
  \hline
  Disiplin tenggat & Lewat $>$1 minggu & Lewat & Tepat waktu & Tepat waktu dan disertai umpan balik \\
  \hline
\end{tabular}
\end{center}

\subsection{Rubrik Laporan dan Kode Akhir (15\%)}

\renewcommand{\arraystretch}{1.35}
\begin{center}\small
\begin{tabular}{|L{9.4cm}|c|}
  \hline
  \thead{\color{ctpBlue} Kriteria} & \thead{\color{ctpBlue}Bobot} \\
  \hline
  \textbf{Kebenaran teknis dan reproduksibilitas} --- eksperimen berjalan dari nol mengikuti \code{README}, \emph{seed} ditetapkan, angka dapat diverifikasi. & 5\% \\
  \hline
  \textbf{Kedalaman analisis dan ablasi} --- pembandingan sebelum/sesudah, analisis jurang lintas bahasa, pengujian \emph{leakage}, dan pembahasan limitasi. & 5\% \\
  \hline
  \textbf{Kualitas penulisan} --- struktur laporan runtut, tabel dan gambar rapi, klaim didukung angka. & 3\% \\
  \hline
  \textbf{Etika dan atribusi} --- pernyataan limitasi, keterwakilan data, dan atribusi CC BY-SA 4.0 lengkap. & 2\% \\
  \hline
  \textbf{Total} & \textbf{15\%} \\
  \hline
\end{tabular}
\end{center}

\subsection{Rubrik Presentasi dan Kontribusi Individu (10\%)}

\renewcommand{\arraystretch}{1.35}
\begin{center}\small
\begin{tabular}{|L{9.4cm}|c|}
  \hline
  \thead{\color{ctpBlue} Kriteria} & \thead{\color{ctpBlue}Bobot} \\
  \hline
  \textbf{Presentasi dan demo} --- alur jelas, demo berjalan, waktu 10 menit dipatuhi. & 4\% \\
  \hline
  \textbf{Tanya jawab} --- penguasaan materi, kemampuan membela keputusan teknis, dan mengakui keterbatasan. & 3\% \\
  \hline
  \textbf{Kontribusi individu} --- log mingguan lengkap, penilaian sebaya, dan ketepatan menjawab pertanyaan tertarget sesuai peran. & 3\% \\
  \hline
  \textbf{Total} & \textbf{10\%} \\
  \hline
\end{tabular}
\end{center}

% =============================================================
\section{Apa yang Dikumpulkan}\label{sec:submit}

Seluruh artefak proyek dikemas ke dalam \textbf{satu berkas arsip} dengan format penamaan:
\begin{center}
  \code{proyek\_akhir\_kelompok\_NN.zip}
\end{center}

dengan \code{NN} diganti oleh nomor kelompok dua digit tanpa spasi, misalnya:
\begin{center}
  \code{proyek\_akhir\_kelompok\_03.zip}
\end{center}

\textbf{Artefak yang wajib disertakan di dalam arsip}:
\begin{borderedsquare}
  \item \code{laporan\_kelompok\_NN.pdf} --- laporan teknis 4--6 halaman.
  \item \code{kode\_kelompok\_NN.zip} atau \code{kode\_kelompok\_NN.ipynb} --- kode
        \emph{end-to-end}. Notebook \textbf{wajib telah dijalankan} sehingga seluruh
        \emph{output} sel tersimpan di dalam berkas.
  \item \code{README.md} --- cara menjalankan kode, daftar dependensi, dan penetapan
        \emph{seed}.
  \item \code{slide\_kelompok\_NN.pdf} --- materi presentasi.
  \item \code{log\_kontribusi\_kelompok\_NN.csv} --- log mingguan per anggota, dengan
        kolom: nama, NIM, peran, minggu, kegiatan, bukti.
  \item \code{hasil\_test\_kelompok\_NN.csv} --- prediksi untuk \emph{test set} privat,
        sesuai format yang ditetapkan dosen.
\end{borderedsquare}

Setiap berkas menggunakan awalan \code{kelompok\_NN} yang sama. Pastikan seluruh berkas
berada tepat di dalam satu arsip \code{.zip} tanpa struktur folder bertingkat yang tidak
perlu.

\begin{callout}[ctpBlue]
\textbf{\color{ctpBlue}Catatan Pengerjaan}
\begin{itemize}
  \item \textbf{Kelompok:} 5 mahasiswa. Nama kelompok dan pembagian peran dikunci pada M0
        dan tidak boleh berubah tanpa persetujuan dosen.
  \item \textbf{Kode yang tidak dapat dijalankan} dinilai sebagai tidak memenuhi syarat
        pada kriteria kebenaran teknis, berapa pun angka yang dilaporkan.
  \item \textbf{Test set} bersifat privat. Satu pengajuan prediksi per kelompok pada M4;
        pengajuan tambahan tidak dilayani.
  \item \textbf{Estimasi waktu:} 2--3 jam per anggota per minggu selama delapan minggu.
\end{itemize}
\end{callout}

\section*{Referensi}
{\small\color{ctpSubtext}
\begin{enumerate}[label={[\arabic*]}, leftmargin=2em]
  \item Cahyawijaya, S., et al. (2025). \emph{Crowdsource, Crawl, or Generate? Creating SEA-VL, a Multicultural Vision-Language Dataset for Southeast Asia}. ACL. --- sumber dataset proyek ini.
  \item Radford, A., et al. (2021). \emph{Learning Transferable Visual Models From Natural Language Supervision}. ICML. --- arsitektur CLIP dan tujuan kontrastif.
  \item Carlsson, F., et al. (2022). \emph{Cross-lingual and Multilingual CLIP}. --- dasar teoretis kesenjangan lintas bahasa pada Soal M2.
  \item Liang, P., Zadeh, A., \& Morency, L.-P. (2023). \emph{Tutorial on Multimodal Machine Learning}. ICML. --- kerangka pilar representasi, penyejajaran, dan fusi.
  \item Thakur, N., et al. (2021). \emph{BEIR: A Heterogeneous Benchmark for Zero-shot Evaluation of Information Retrieval Models}. NeurIPS Datasets and Benchmarks. --- rujukan protokol evaluasi retrieval.
  \item Pineau, J., et al. (2021). \emph{Improving Reproducibility in Machine Learning Research}. JMLR. --- dasar syarat reproduksibilitas kode.
  \item Gebru, T., et al. (2021). \emph{Datasheets for Datasets}. Communications of the ACM. --- kerangka dokumentasi dataset dan sumber pemeriksaan bias.
  \item Bender, E. M., et al. (2021). \emph{On the Dangers of Stochastic Parrots: Can Language Models Be Too Big?}. FAccT. --- dasar diskusi keterwakilan dan dampak sosial pada M3.
\end{enumerate}}

\end{document}
